\section{Data}
\label{sec:data}

The analysis uses the harmonised Households Below Average Income (HBAI) extract of the Family Resources Survey (UKDA-5828), a repeated cross-section --- a different sample of households is drawn each year, not the same households tracked over time --- covering financial years ending (FYE) 2017--2024, in 2023/24 prices, and excluding FYE2021 because of Covid-19 fieldwork disruption. The sample is restricted to children aged 16 or under living in England (the control group) or Scotland (the treatment group); Wales and Northern Ireland are excluded from the main design and used instead as falsification checks (Section~\ref{sec:results}). Table~\ref{tab:summary_sample_composition} reports the resulting sample composition by survey year and country: England contributes the large majority of observations in every year, while Scotland's smaller population share is reflected in correspondingly smaller cell counts, ranging from 537 (FYE2024) to 1,190 (FYE2017) children.

\input{../tables/latex/summary_sample_composition.tex}

The Scottish Child Payment (SCP) was introduced for under-6s in benefit-receiving families from 15 February 2021 (\pounds10/week, raised to \pounds20/week from April 2022) and expanded to all under-16s at \pounds25/week from 14 November 2022. Following Nesom, Stewart and Tominey (2025) and Andersen et al.\ (2025), the headline design treats the November 2022 expansion as the single treatment date, with the post-treatment indicator refined to the exact 14 November cutoff within FYE2023 where household interview dates are available from the underlying FRS file; the earlier, age-limited 2021 rollout is not used to define treatment timing in the headline specification, since only a small under-6 subgroup is affected and statistical power to detect a separate age-cohort effect is limited (Section~\ref{sec:results} nonetheless reports an age-cohort robustness check where feasible).

The primary outcome is the official DWP child material deprivation (MDCH) flag, a binary indicator equal to one if a child's household is classified as materially deprived under DWP's scoring methodology. Because the underlying MDCH questionnaire was redesigned at FYE2024 (moving from a 21-item, threshold-based design to a revised 11+11-item, simple-count design), ten individual deprivation items common to both designs are also retained and recoded to 0/1 (deprived/not deprived). From these ten items, three composite measures are constructed on a single, unrevised definition throughout the panel: \emph{any deprivation} (at least one item lacking), \emph{deprivation count} (the mean number of items lacking), and \emph{severe deprivation} (three or more items lacking, this dissertation's own approximate measure rather than an official DWP statistic). A secondary food-insecurity outcome, available only from FYE2020, is used for comparability with Stewart et al.\ (2025) and as an additional pre-trend diagnostic. The deprivation module is administered on an identical basis across Great Britain, with no country-specific question wording or routing in the FRS instrument, so the measure is treated as consistently defined across the Scotland--England comparison.

Table~\ref{tab:summary_outcomes} summarises the four headline outcomes by country and period. Both countries show a rise in any-deprivation and severe-deprivation rates over time, but the rise is markedly larger in England (any deprivation: 18.9\% to 26.5\%) than in Scotland (14.8\% to 15.9\%) --- the raw divergence that motivates the difference-in-differences design below. Food insecurity, by contrast, rises by a similar margin in England but stays essentially flat in Scotland.

\input{../tables/latex/summary_outcomes.tex}

Table~\ref{tab:summary_background} reports weighted background characteristics for England and Scotland pooled across all years, replicating Table 2 of Stewart et al.\ (2025, CASE). Scotland has a smaller share of large (3+ child) families, a similar share of lone-parent households, a slightly younger and more female-headed household-head profile, and a substantially larger White household share (91.8\% vs.\ 76.6\%) than England. Table~\ref{tab:summary_a1_by_period} decomposes these same characteristics by pre-/post-SCP period within each country (split at the November 2022 cutoff), replicating Stewart et al.'s Table A1; the pre/post composition of each country's sample is broadly stable, with the exception of a rise in the disabled-household share in both countries and a fall in Scotland's household size in the post period, consistent with the identifying assumption that observable composition is not shifting sharply around the treatment date in a way that would itself explain a change in measured deprivation.

\input{../tables/latex/summary_background.tex}
\input{../tables/latex/summary_a1_by_period.tex}

Two covariate sets are used throughout. A lean, six-variable set replicates the CASE evaluation's specification exactly (Stewart et al., 2025): a dummy for household head aged under 25, female household head, a five-category ethnicity factor, disability status of the household, lone-parent household, and large family (three or more dependent children). Household income, benefit-receipt flags, and housing tenure are deliberately excluded from this set: none is part of CASE's published specification, and including them carries a genuine mediator/bad-control risk here, since income responds mechanically to receiving SCP and benefit receipt was itself the SCP eligibility gate before the November 2022 expansion. A wide, approximately 35-variable set --- continuous income, housing and benefit-unit amounts, family composition, and all nine benefit-receipt flags, one-hot encoded where multi-level --- is used only in the double/debiased machine learning (DML) analysis (Section~\ref{sec:results}), to test sensitivity to covariate breadth once the estimator no longer requires a hand-specified linear functional form.

All estimates are weighted by the dependent-child survey weight (\texttt{GS\_INDCH}); standard errors are clustered at household level throughout (Section~\ref{sec:methodology}). The analytic sample is approximately 47,000 child-year observations for the item-based outcomes and 51,000 for the official MDCH flag.

\section{Methodology}
\label{sec:methodology}

The empirical analysis proceeds in five completed stages, following a design agreed with the project supervisor: (1) baseline covariate-adjusted OLS difference-in-differences (DiD) on the single official outcome; (2) the same design re-estimated jointly across all ten deprivation items in one household-clustered model; (3) a doubly robust machine-learning re-estimation of the baseline, testing sensitivity to functional form and covariate breadth; (4) the same machine-learning approach applied separately to each of the ten items; and (5) a jointly estimated, vector-outcome version of stage (4). A sixth stage --- an age-cohort staggered-timing extension reflecting SCP's true two-phase rollout --- is explored only as a background robustness check, time permitting, since the primary cross-fitted estimator used here is not designed for staggered adoption timing.

\subsection{Baseline difference-in-differences}
\label{sec:baseline-did}

The estimating equation for the single-outcome baseline is
\begin{equation}
Y_{it} = \alpha + \beta\,\text{Treated}_{i} + \delta\,(\text{Treated}_{i} \times \text{Post}_{t}) + X_{it}'\theta + \lambda_{t} + \varepsilon_{it},
\label{eq:baseline-did}
\end{equation}
fitted by weighted least squares with household-clustered standard errors, where $\lambda_{t}$ are survey-year fixed effects and $\delta$, the coefficient on the Treated$\times$Post interaction, is the DiD estimate of the average treatment effect on the treated (ATT). This is numerically equivalent to the standard $2\times2$ estimand once covariates are dropped. Three variants are reported: \emph{Simple} (no covariates beyond the DiD terms and year fixed effects), \emph{Adjusted} (the six CASE controls described in Section~\ref{sec:data}, the genuine Stewart et al.\ replication specification), and \emph{Extended} (Adjusted plus the child's own age, a labelled robustness variant, since several deprivation items are age-differentiated in what they capture). An item-by-item version of Equation~\ref{eq:baseline-did}, with Benjamini--Hochberg (1995) false-discovery-rate correction applied across the ten resulting tests, is also estimated, alongside an exact replication of Stewart et al.'s (2025) own Table 3 specification (which includes an explicit \texttt{Post} main-effect term rather than absorbing it into year fixed effects, and reports the full covariate vector).

\subsection{Item-stacked difference-in-differences}
\label{sec:item-stacked}

Fitting the ten deprivation items as ten independent regressions treats each item's error term as unrelated to the others, but the same child answers all ten deprivation questions, so their responses are correlated within household --- a Moulton (1990) problem. This stage reshapes the ten items into long format (one row per child-year-item) and fits a single joint model with item and year fixed effects, clustering standard errors at household level so that they correctly reflect this within-child correlation. Both a pooled specification (one average effect across all ten items) and an item-interacted specification (one coefficient per item, with a joint Wald test of the null that all ten item-specific effects are zero) are estimated, in both Simple and Adjusted forms using the same covariates as Section~\ref{sec:baseline-did}.

\subsection{Identification: parallel trends}
\label{sec:parallel-trends}

Parallel pre-treatment trends are tested via Scotland$\times$time interactions, replicating the CASE evaluation's own diagnostic tables: a covariate balance table (Scotland$\times$year interactions for each background characteristic) and an outcome event-study table (Scotland$\times$year interactions for each outcome and item), both tested jointly via a Wald test with Benjamini--Hochberg correction applied across outcomes. An initial pass using region-clustered standard errors (clustering at Government Office Region, of which there are only around ten, with Scotland the sole treated region) flagged almost every covariate and outcome as a possible pre-trend, including purely demographic covariates with no plausible SCP channel --- a signal that the test itself was miscalibrated rather than that trends were genuinely diverging. A placebo test, re-assigning "treated" status to each English region in turn, confirmed this: the large majority of placebo regions "failed" the identical pre-trend test with no policy applied, and Scotland was not an outlier within that distribution. This is the well-documented few-clusters, single-treated-cluster problem in cluster-robust inference (Cameron and Miller, 2015; MacKinnon and Webb, 2018): with only around ten clusters and Scotland the only treated one, cluster-robust joint Wald tests are anti-conservative. Standard errors were instead re-clustered at household level throughout the pipeline, which also resolved the miscalibration and is independently the correct level given the within-child correlation described in Section~\ref{sec:item-stacked}.

At annual granularity, and under Benjamini--Hochberg correction, every outcome and covariate in both diagnostic tables reads as consistent with parallel trends. A separate, finer-grained check replicates the CASE paper's own event-study specification at quarterly granularity, using the paper's rolling 13-week windows anchored on the SCP rollout date (14 February/May/August/November each year) rather than annual survey-year dummies. Table~\ref{tab:parallel_trends_quarterly_wald} reports the resulting joint Wald tests, Benjamini--Hochberg corrected across all 15 outcomes and items tested. At this finer granularity, three series remain flagged even after correction: any deprivation (mdch\_any), Bed/bedroom, and Outdoor play area. This does not overturn the annual-level conclusion for the headline official MDCH flag or the deprivation-count and severe-deprivation composites, all of which remain clean at quarterly granularity, but it qualifies the strength of the identifying assumption specifically for the any-deprivation composite and two individual items, and is reported here as an explicit limitation rather than resolved.

\input{../tables/latex/parallel_trends_quarterly_wald.tex}

\subsection{Double/debiased machine learning difference-in-differences}
\label{sec:dml}

Stage 3 asks whether the Section~\ref{sec:baseline-did} result depends on OLS's linear functional form or on its narrow six-variable control set, rather than on the treatment itself. The Adjusted specification in Equation~\ref{eq:baseline-did} removes confounding by assuming a specific, linear relationship between the six CASE covariates and deprivation; if that relationship is in fact nonlinear, or involves interactions between covariates, a linear model is misspecified and the DiD coefficient absorbs some of that misspecification as bias. Substituting a flexible machine-learning model for the covariate adjustment creates a new problem, since regularised ML methods bias whatever moment condition is built on top of them. Chernozhukov et al.\ (2018) show this contamination can be removed by combining a Neyman-orthogonal moment condition with cross-fitting, so that the ML model's own overfitting cannot leak into the causal estimate.

The estimator used is Zimmert's (2020) efficient-score, doubly robust DiD estimator, implemented via \texttt{causalweight::didDML()}. For each of the four treatment$\times$time cells, the estimator fits a cross-fitted outcome-regression model and a cell-propensity model, combined into an augmented score that is consistent if either the outcome-regression or the propensity model is correctly specified in each cell, not necessarily both, and Neyman-orthogonal, so that cross-fitted ML estimates in place of a correctly specified parametric model do not bias the ATT estimate at first order. An earlier candidate construction, Chang's (2020) repeated-cross-section estimator, was implemented first but discarded after it returned an implausibly large, tightly estimated effect consistent with propensity-model collapse; Zimmert's construction was adopted in its place.

Two nuisance-model methods are compared: lasso, a linear method that (like OLS) cannot represent nonlinearities or interactions, and random forest, which can. Two covariate sets are also compared: the lean six CASE controls (isolating the effect of switching estimators alone) and the wide $\sim$35-variable set (testing whether covariate breadth matters once the estimator no longer requires a hand-specified linear form); this wide-set comparison is only feasible because of cross-fitting, since an earlier attempt to fit the same wide set within a group-time ATT (Callaway--Sant'Anna-style) framework produced singular, rank-deficient models. Table~\ref{tab:stage3_dml_composite} reports the resulting ATET estimates for the official MDCH flag under all combinations, plus a lasso/random-forest ensemble estimated for the lean specification only, since the ensemble did not complete within a feasible runtime for the wide specification.

\input{../tables/latex/stage3_dml_composite.tex}

\subsection{Item-level and stacked machine-learning DiD}
\label{sec:item-dml}

The Section~\ref{sec:dml} estimator is repeated separately for each of the ten deprivation items, using the wide covariate set and lasso as the primary nuisance-model method (random forest added for a second pass), with Benjamini--Hochberg correction applied across items. A fifth stage extends this to a single, jointly estimated model across all ten items using a shared complete-case sample and a multivariate outcome-regression learner, since no existing DML-DiD implementation accepts a vector-valued outcome directly; a validation step reconstructing a single item's result by hand, using the same underlying nuisance-fitting machinery, confirmed the joint reimplementation is faithful before it was used for inference. Both the per-item and the jointly estimated results are reported in Section~\ref{sec:results}, once the item-label correction described there has been propagated through to the underlying result files.

\subsection{Robustness}
\label{sec:robustness}

Three further checks are reported alongside the main results. First, a specification excluding FYE2022 (the partial-rollout year, in which under-6 Scottish children were already receiving SCP for part of the year) from the pre-period is estimated alongside the Section~\ref{sec:baseline-did} baseline. Second, an official-MDCH-flag measurement-consistency check confirms the flag is computed on the same basis in Scotland and England, with no region-specific necessities weighted differently between the two nations. Third, a falsification design re-estimates the headline specification with Wales and Northern Ireland in turn as the treated unit against England as control, using SCP's actual rollout date as a pseudo-treatment date; since neither nation received SCP, a null coefficient is the reassuring result, supporting attribution of the Scotland--England divergence to the policy itself rather than to a UK-wide shock coincident with SCP's timing.
